% ============================================================================
% Fichier  : partie1/P01_C01_securite_information.tex
% Titre    : Qu'est-ce que la Sécurité de l'Information
% Auteur   : MINKA MI NGUIDJOÏ Thierry Emmanuel
% ============================================================================
% Sources : ISO/IEC 27000:2018 ; ISO/IEC 27001:2022. Mairie d'Eséka : cas fictif.
% ============================================================================

\chapter{Qu'est-ce que la Sécurité de l'Information}
\label{chap:cia}
\chaptermark{Sécurité de l'information}

\epigraph{%
    « Protéger l'information, ce n'est pas la cacher~:\\
    c'est la rendre sûre, exacte et disponible à qui de droit. »%
}{Principe du cours}

\niveauChapitre{Fondamental}{Le prologue. Aucune connaissance technique
préalable.}

\begin{boxobjectifs}
\begin{itemize}
    \item Définir la sécurité de l'information et la distinguer de la sécurité
          informatique.
    \item Énoncer la triade \sigle{DIC} et l'illustrer sur des actifs réels.
    \item Ajouter les critères d'authenticité, de non-répudiation et de
          traçabilité.
    \item Identifier l'information comme un actif à valeur, quel que soit son
          support.
    \item Relier chaque critère à une conséquence concrète pour la mairie.
\end{itemize}
\end{boxobjectifs}

\bigskip

% ============================================================================
\section*{Question directrice : que protège-t-on, au juste, quand on « sécurise » ?}
% ============================================================================

La clé \texttt{USB} a chiffré des fichiers~: l'information est devenue
indisponible. Mais si elle avait été copiée et revendue, le problème aurait
été autre. Ce chapitre nomme précisément ce que l'on protège et selon quels
critères.

\medskip
\textbf{Plan du chapitre.}
\begin{enumerate}
    \item L'information, un actif (section~\ref{sec:c01-actif}).
    \item La triade \sigle{DIC} (section~\ref{sec:c01-dic}).
    \item Les critères complémentaires (section~\ref{sec:c01-criteres}).
    \item Sécurité de l'information ou sécurité informatique (section~\ref{sec:c01-perimetre}).
    \item Lire un incident par ses critères (section~\ref{sec:c01-lecture}).
\end{enumerate}

\separateur

% ============================================================================
\section{L'Information, un Actif}
\label{sec:c01-actif}
% ============================================================================

\begin{boxdef}
\termecle{Information}~: donnée qui a un sens et une valeur pour
l'organisation, quel que soit son support (fichier, papier, parole, savoir
d'un agent). \termecle{Actif}~: tout élément ayant de la valeur et qu'il faut
donc protéger.
\end{boxdef}

On sécurise une information parce qu'elle a de la valeur~: les actes de
naissance font foi, le fichier des contribuables conditionne les recettes, la
messagerie engage la commune. La valeur n'est pas dans le disque dur~: elle est
dans l'information. Protéger le support sans penser l'information, c'est
protéger le coffre en oubliant ce qu'il contient.

\begin{boxobjet}{Le registre d'état civil numérisé}
\textbf{Constatation.} Le registre numérisé des naissances est consultable par
tous les agents du service, sans mot de passe individuel, et n'est copié nulle
part ailleurs.

\medskip
\textbf{Analyse.} C'est un actif de valeur majeure (il fait foi) mais sans
protection~: ni confidentialité (accès non maîtrisé), ni disponibilité (copie
unique). Sa valeur et sa vulnérabilité en font une cible prioritaire.
\end{boxobjet}

% ============================================================================
\section{La Triade DIC}
\label{sec:c01-dic}
% ============================================================================

\begin{boxdef}
La \termecle{triade \sigle{DIC}} définit trois critères de sécurité~:
\begin{itemize}
    \item \termecle{Disponibilité}~: l'information est accessible à qui de
          droit quand il en a besoin.
    \item \termecle{Intégrité}~: l'information n'est ni altérée ni détruite de
          façon non autorisée.
    \item \termecle{Confidentialité}~: l'information n'est accessible qu'aux
          personnes autorisées.
\end{itemize}
En anglais~: \emph{CIA} (Confidentiality, Integrity, Availability).
\end{boxdef}

Les trois critères se tiennent et s'arbitrent. Verrouiller un registre
(confidentialité) sans en garder de copie met en péril sa disponibilité. La
sécurité consiste à tenir les trois, à hauteur de la valeur de l'actif.

\begin{tableaucours}{L{3.0cm}L{4.6cm}L{4.6cm}}{%
    \tableauHeader{Critère} &
    \tableauHeader{Atteinte à Eséka} &
    \tableauHeader{Conséquence}
}
Disponibilité & Registre chiffré, illisible & Service à l'arrêt, usagers bloqués \\
Intégrité & Date de naissance modifiée & Acte faux, responsabilité engagée \\
Confidentialité & Fichier contribuables copié & Vie privée exposée, loi 2024/017 \\
\end{tableaucours}
\vspace{-0.8em}
\begin{center}{\small\itshape Tableau~1.1. La triade DIC sur les actifs de la mairie}\end{center}

% ============================================================================
\section{Les Critères Complémentaires}
\label{sec:c01-criteres}
% ============================================================================

À la triade s'ajoutent, selon les besoins, trois critères~:

\begin{itemize}
    \item \termecle{Authenticité}~: l'origine de l'information est vérifiée
          (l'acte vient bien de la mairie).
    \item \termecle{Non-répudiation}~: on ne peut nier avoir effectué une
          action (l'agent qui a validé un acte ne peut le contester).
    \item \termecle{Traçabilité}~: les actions sont enregistrées et
          rattribuables.
\end{itemize}

Ces critères soutiennent la confiance~: un acte authentique, non répudiable et
tracé a une valeur juridique qu'un simple fichier n'a pas.

% ============================================================================
\section{Sécurité de l'Information ou Sécurité Informatique}
\label{sec:c01-perimetre}
% ============================================================================

\begin{boxdef}
La \termecle{sécurité de l'information} protège l'information sous toutes ses
formes (numérique, papier, orale). La \termecle{sécurité informatique} ne
protège que les systèmes qui la traitent. La première englobe la seconde.
\end{boxdef}

La distinction n'est pas scolaire. À la mairie, un acte lu à voix haute au
guichet, un dossier papier oublié sur un bureau ou une confidence échappent à
toute mesure informatique mais relèvent de la sécurité de l'information. Un
\sigle{SMSI} couvre l'information, pas seulement les ordinateurs.

\begin{boxmetier}{RSSI}
\textbf{Penser actif, pas machine.} Le \sigle{RSSI} d'Eséka commence par
lister les informations de valeur et leur support, avant de parler pare-feu.
On protège ce qui a de la valeur, pas ce qui a une prise électrique.
\end{boxmetier}

% ============================================================================
\section{Lire un Incident par ses Critères}
\label{sec:c01-lecture}
% ============================================================================

Nommer les critères donne une grille de lecture. Tout incident s'analyse en
demandant quel critère a été atteint, sur quel actif, et avec quelle
conséquence.

\begin{boxprocedure}
\textbf{Constatation.} Le 13~janvier, les fichiers du service d'état civil
sont illisibles.

\medskip
\textbf{Analyse} (confiance élevée). Critère atteint~: \emph{disponibilité}.
Actif~: registres et tableurs du service. L'intégrité et la confidentialité
ne sont pas, à ce stade, mises en cause~: distinguer les trois évite de
conclure trop vite à un vol de données.
\end{boxprocedure}

\separateur

\begin{boxresume}
\textbf{Ce qu'il faut retenir de ce chapitre~:}
\begin{itemize}
    \item L'\textbf{information} est un actif à valeur~; on la protège pour
          cette valeur.
    \item La \textbf{triade DIC} (disponibilité, intégrité, confidentialité)
          définit ce que « sécuriser » veut dire.
    \item \textbf{Authenticité, non-répudiation, traçabilité} soutiennent la
          confiance.
    \item La \textbf{sécurité de l'information} dépasse la sécurité
          informatique.
    \item Tout incident se lit par le \textbf{critère atteint} sur un actif.
\end{itemize}
\end{boxresume}

\begin{boxinfo}
\textbf{Points de contrôle de l'élève.} Je sais~: définir information et
actif~; énoncer la triade \sigle{DIC}~; citer les trois critères
complémentaires~; distinguer sécurité de l'information et sécurité
informatique~; analyser un incident par ses critères.
\end{boxinfo}

\bigskip

\begin{boxexo}[title={\ding{45}\quad Exercice~1.1. Trois critères \niveaubase}]
Pour chacun des cas suivants, dites quel critère \sigle{DIC} est atteint~:
registre chiffré~; montant d'une taxe modifié~; liste d'électeurs copiée.
\end{boxexo}

\begin{boxexo}[title={\ding{45}\quad Exercice~1.2. Actif ou support \niveaubase}]
Parmi~: le disque dur, l'acte de naissance, l'ordinateur, le savoir de l'agent,
lesquels sont des actifs informationnels~? Justifiez.
\end{boxexo}

\begin{boxexo}[title={\ding{45}\quad Exercice~1.3. Information ou informatique \niveaubase}]
Donnez un exemple, à la mairie, d'atteinte à la sécurité de l'information qui
ne passe par aucun ordinateur.
\end{boxexo}

\begin{boxexo}[title={\ding{45}\quad Exercice~1.4. Les compléments \niveaumoyen}]
Expliquez, sur l'exemple d'un acte de naissance, l'authenticité, la
non-répudiation et la traçabilité.
\end{boxexo}

\begin{boxexo}[title={\ding{45}\quad Exercice~1.5. Arbitrage \niveaumoyen}]
Verrouiller le registre sans copie renforce un critère et en menace un autre.
Lesquels~? Comment tenir les deux~?
\end{boxexo}

\begin{boxexo}[title={\ding{45}\quad Exercice~1.6. Valeur d'un actif \niveaumoyen}]
Classez trois actifs de la mairie par valeur décroissante et justifiez le
classement.
\end{boxexo}

\begin{boxexo}[title={\ding{45}\quad Exercice~1.7. Grille de lecture \niveaumoyen}]
Appliquez la grille « critère / actif / conséquence » au vol du fichier des
contribuables.
\end{boxexo}

\begin{boxexo}[title={\ding{45}\quad Exercice~1.8. Confidentialité et loi \niveaumoyen}]
En quoi une atteinte à la confidentialité du fichier des contribuables
engage-t-elle la mairie au regard de la loi 2024/017~?
\end{boxexo}

\begin{boxexo}[title={\ding{45}\quad Exercice~1.9. Cartographie \niveauexpert}]
Dressez un tableau des cinq principaux actifs informationnels de la mairie
d'Eséka, avec leur support et les critères prioritaires pour chacun.
\end{boxexo}

\begin{boxexo}[title={\ding{45}\quad Exercice~1.10. Note de cadrage \niveauexpert}]
Rédigez une note d'une page au maire définissant « sécurité de l'information »
pour la mairie et justifiant pourquoi elle dépasse l'achat d'un antivirus.
\end{boxexo}

\bigskip

\begin{boxinfo}
\textbf{Pour aller plus loin.}
\begin{itemize}
    \item ISO/IEC 27000~\cite{iso27000}~; ISO/IEC 27001~\cite{iso27001}.
\end{itemize}
\end{boxinfo}

\begin{boxinfo}
\textbf{Transition.} Savoir ce que l'on protège ne dit pas encore pourquoi ni
pour qui~: le Chapitre~\ref{chap:enjeux} relie le risque à la décision.
\end{boxinfo}
